% Blazor Server section

% 1 – What is Blazor Server?
\QuestionSlide[\CategoryBadge[BlazorColor!20]{Blazor Server}]<1>{What is Blazor Server?}

\begin{frame}[fragile]
  \frametitle{\AnswerTitle[\CategoryBadge[BlazorColor!20]{Blazor Server}]{What is Blazor Server?}}

  {\footnotesize
  \textbf{Blazor Server} runs C\# UI components on the server. The browser renders a thin HTML shell; user events (clicks, input) travel over a \textbf{SignalR WebSocket} to the server, which updates a virtual DOM diff and pushes only the changes back. This keeps all logic and data server-side but requires persistent connections.
  }
\end{frame}

% 2 – Components
\QuestionSlide[\CategoryBadge[BlazorColor!20]{Blazor Server}]<1>{What is a Blazor component?}

\begin{frame}[fragile]
  \frametitle{\AnswerTitle[\CategoryBadge[BlazorColor!20]{Blazor Server}]{What is a Blazor component?}}

  {\footnotesize
  A \textbf{component} is a self-contained unit of UI defined in a \texttt{.razor} file. It combines HTML markup, C\# logic in \texttt{@code\{\}}, and optionally CSS isolation. Components receive data through \texttt{[Parameter]} properties and communicate upward via \texttt{EventCallback}.
  }

  \begin{minted}[fontsize=\scriptsize]{html}
<!-- Counter.razor -->
<h3>Count: @count</h3>
<button @onclick="Increment">+1</button>

@code {
    [Parameter] public int InitialValue { get; set; }
    private int count;
    protected override void OnInitialized() => count = InitialValue;
    private void Increment() => count++;
}
  \end{minted}
\end{frame}

% 3 – Component Lifecycle
\QuestionSlide[\CategoryBadge[BlazorColor!20]{Blazor Server}]<2>{What is the Blazor component lifecycle?}

\begin{frame}[fragile]
  \frametitle{\AnswerTitle[\CategoryBadge[BlazorColor!20]{Blazor Server}]{What is the Blazor component lifecycle?}}

  {\footnotesize
  Key lifecycle methods (sync + async variants):
  \begin{enumerate}
    \item \texttt{SetParametersAsync} – parameters set from parent.
    \item \texttt{OnInitialized[Async]} – one-time init; fetch initial data here.
    \item \texttt{OnParametersSet[Async]} – called when parameters change.
    \item \texttt{OnAfterRender[Async]} – DOM is ready; safe for JS interop.
    \item \texttt{Dispose} / \texttt{DisposeAsync} – clean up timers, subscriptions.
  \end{enumerate}
  }

  \begin{minted}[fontsize=\scriptsize]{csharp}
@implements IAsyncDisposable

@code {
    private List<Product> products = [];

    protected override async Task OnInitializedAsync()
        => products = await ProductService.GetAllAsync();

    public async ValueTask DisposeAsync()
        => await hubConnection.DisposeAsync();
}
  \end{minted}
\end{frame}

% 4 – Data Binding
\QuestionSlide[\CategoryBadge[BlazorColor!20]{Blazor Server}]<1>{What is data binding in Blazor?}

\begin{frame}[fragile]
  \frametitle{\AnswerTitle[\CategoryBadge[BlazorColor!20]{Blazor Server}]{What is data binding in Blazor?}}

  {\footnotesize
  \begin{itemize}
    \item \textbf{One-way} – \texttt{@variable} renders a value; updates do not flow back automatically.
    \item \textbf{Two-way} – \texttt{@bind="property"} synchronises between the UI element and C\# field on the \texttt{oninput} or \texttt{onchange} event.
    \item \textbf{Event-driven} – \texttt{@onclick}, \texttt{@oninput} call C\# handlers that update state and trigger re-render.
  \end{itemize}
  }

  \begin{minted}[fontsize=\scriptsize]{html}
<!-- two-way binding to a text field -->
<input @bind="searchTerm" @bind:event="oninput" />
<p>Searching: @searchTerm</p>

@code {
    private string searchTerm = "";
}
  \end{minted}
\end{frame}

% 5 – State Management
\QuestionSlide[\CategoryBadge[BlazorColor!20]{Blazor Server}]<2>{How do you manage state in Blazor Server?}

\begin{frame}[fragile]
  \frametitle{\AnswerTitle[\CategoryBadge[BlazorColor!20]{Blazor Server}]{How do you manage state in Blazor Server?}}

  {\footnotesize
  Options from simplest to most powerful:
  \begin{itemize}
    \item \textbf{Component fields} – local, lost on reconnect.
    \item \textbf{Scoped service} (injected) – lives per SignalR circuit; shared between components in the same browser tab.
    \item \textbf{CascadingValue} – pass state down the component tree without prop-drilling.
    \item \textbf{Browser storage} (via JS interop) – survives page refresh.
    \item \textbf{Fluxor / Blazor-State} – Redux-style store for complex apps.
  \end{itemize}
  }
\end{frame}

% 6 – JS Interop
\QuestionSlide[\CategoryBadge[BlazorColor!20]{Blazor Server}]<2>{What is JavaScript Interop in Blazor?}

\begin{frame}[fragile]
  \frametitle{\AnswerTitle[\CategoryBadge[BlazorColor!20]{Blazor Server}]{What is JavaScript Interop?}}

  {\footnotesize
  \textbf{JS Interop} lets Blazor call JavaScript functions and JS call .NET methods. Use \texttt{IJSRuntime.InvokeAsync<T>} for calls that return values; \texttt{InvokeVoidAsync} when no return is needed. Only safe inside \texttt{OnAfterRenderAsync} when the DOM exists.
  }

  \begin{minted}[fontsize=\scriptsize]{csharp}
@inject IJSRuntime JS

@code {
    protected override async Task OnAfterRenderAsync(bool firstRender)
    {
        if (firstRender)
            await JS.InvokeVoidAsync("initChart", chartRef);
    }

    private async Task CopyToClipboard(string text)
        => await JS.InvokeVoidAsync("navigator.clipboard.writeText", text);
}
  \end{minted}
\end{frame}

% 7 – EditForm & Validation
\QuestionSlide[\CategoryBadge[BlazorColor!20]{Blazor Server}]<2>{How do you handle forms and validation in Blazor?}

\begin{frame}[fragile]
  \frametitle{\AnswerTitle[\CategoryBadge[BlazorColor!20]{Blazor Server}]{How do you handle forms and validation in Blazor?}}

  {\footnotesize
  \texttt{EditForm} wraps an \texttt{EditContext} bound to a model. \texttt{DataAnnotationsValidator} reads \texttt{[Required]}, \texttt{[Range]}, etc. \texttt{ValidationSummary} and \texttt{ValidationMessage<T>} display errors. \texttt{OnValidSubmit} fires only when validation passes.
  }

  \begin{minted}[fontsize=\scriptsize]{html}
<EditForm Model="order" OnValidSubmit="PlaceOrder">
    <DataAnnotationsValidator />
    <ValidationSummary />

    <InputText @bind-Value="order.CustomerName" />
    <ValidationMessage For="@(() => order.CustomerName)" />

    <button type="submit">Place Order</button>
</EditForm>
  \end{minted}
\end{frame}

% 8 – Render Modes (.NET 8)
\QuestionSlide[\CategoryBadge[BlazorColor!20]{Blazor Server}]<3>{What are Blazor render modes in .NET 8?}

\begin{frame}[fragile]
  \frametitle{\AnswerTitle[\CategoryBadge[BlazorColor!20]{Blazor Server}]{What are Blazor render modes in .NET 8?}}

  {\footnotesize
  .NET 8 unified all Blazor hosting models under one project type with \textbf{per-component render modes}:
  \begin{itemize}
    \item \texttt{@rendermode InteractiveServer} – SignalR circuit, server-side logic.
    \item \texttt{@rendermode InteractiveWebAssembly} – runs in browser WASM.
    \item \texttt{@rendermode InteractiveAuto} – WASM on return visits, Server first load.
    \item \textbf{Static SSR} (default) – HTML only, no interactivity.
  \end{itemize}
  }

  \begin{minted}[fontsize=\scriptsize]{html}
<!-- Per-page render mode -->
@rendermode InteractiveServer

<!-- Per-component in parent -->
<Counter @rendermode="InteractiveWebAssembly" />
  \end{minted}
\end{frame}
